\documentclass{ctexbook}

\usepackage{anyfontsize}
\usepackage{amsmath}
\usepackage{graphicx,subfigure}
\usepackage{bm}
\usepackage{geometry}
\geometry{top=3cm, left=2.5cm, right=2.5cm, bottom=3cm}
\usepackage[shortlabels]{enumitem}
\usepackage{caption}
\captionsetup{font=Large}


\begin{document}\Large

    \title{\textbf{实验五\ \ 测量薄透镜的焦距}}
    \author{\Large 1900011604 张植竣}
    \date{\Large 2020年9月30日}

    \maketitle

    \newpage
    \section*{\zihao{3} 1\ \ 实验步骤和测量数据处理}
    \bigskip

    \subsection*{\zihao{-3} 1.1\ \ 位移法测量薄凸透镜焦距}

    \noindent 1. 粗测凸透镜焦距$f$，取物与屏（像）间距离$A > 4f$，记录物的位置$x_1$和屏的位置$x_2$；

    \noindent 2. 移动凸透镜，使物在屏上成大像，记录此时凸透镜光心位置$x_3$；

    \noindent 3. 移动凸透镜，使物在屏上成小像，记录此时凸透镜光心位置$x_4$；

    \noindent 4. 计算$A = |x_2-x_1|$，$l = |x_4-x_3|$，得到焦距：

    \[f = \frac{A^2 - l^2}{4A}\]

    \noindent \textbf{测量数据如下：}
    \begin{table}[ht!]
        \begin{center}
            \begin{tabular}{|c|c|c|c|c|c|c|}
                \hline
                $x_1$/cm & $x_2$/cm & $x_3$/cm & $x_4$/cm & $A$/cm & $l$/cm & $f$/cm \\
                \hline
                25.00 & 90.00 & 48.16 & 67.26 & 65.00 & 19.10 & 14.85 \\
                \hline
                25.00 & 100.00 & 45.69 & 79.52 & 75.00 & 19.10 & 14.94 \\
                \hline
                25.00 & 110.00 & 44.37 & 90.78 & 85.00 & 19.10 & 14.91 \\
                \hline
            \end{tabular}
            \caption{位移法测量薄凸透镜焦距}
        \end{center}
    \end{table}

	\noindent \textbf{测量得到焦距平均值为：}$\bm{\bar{f} =} \textbf{14.90\,cm}$

    \bigskip
    \subsection*{\zihao{-3} 1.2\ \ 自准直法测量薄凸透镜焦距}

    \noindent 1. 粗测凸透镜焦距$f$，取物与凸透镜光心间距离$\Delta x \approx f$，调整平面镜使之离凸透镜光心较近；

    \noindent 2. 移动凸透镜，使物在物屏上成等大倒立实像，记录此时物的位置$x_1$和凸透镜光心位置$x_2$，记录平面镜位置$x_3$；

    \noindent 3. 计算$\Delta x = |x_2 - x_1|$，得到焦距：
    \[f = \Delta x\]

    \noindent \textbf{测量数据如下：}
    \begin{table}[ht!]
        \begin{center}
            \begin{tabular}{|c|c|c|c|c|}
                \hline
                $x_1$/cm & $x_2$/cm & $x_3$/cm & $\Delta x$/cm & $f$/cm \\
                \hline
                25.00 & 39.83 & 47.50 & 14.83 & 14.83 \\
                \hline
            \end{tabular}
            \caption{自准直法测量薄凸透镜焦距}
        \end{center}
    \end{table}

    \noindent \textbf{测量得到焦距为：}$\bm{f =} \textbf{14.83\,cm}$

    \bigskip
    \subsection*{\zihao{-3} 1.3\ \ 物像距法测量薄凹透镜焦距}

    \noindent 1. 粗测凸透镜焦距$f$，取物与屏（像）间距离$A>4f$，固定屏的位置，记录屏的位置$x_1$；

    \noindent 2. 移动凸透镜，使物在屏上成小像，固定凸透镜的位置，该像将成为凹透镜的虚物，虚物的位置即为$x_1$；

    \noindent 3. 将凹透镜放置在凸透镜和屏之间，移动凹透镜和屏，使虚物在屏上成放大倒立实像，记录此时凹透镜光心位置$x_2$与屏（实像）的位置$x_3$；

    \noindent 4. 计算$u = -|x_2 - x_1|$，$v = |x_3-x_2|$，得到焦距：
    \[f = \frac{uv}{u+v}\]

    \noindent \textbf{测量数据如下：}
    \begin{table}[ht!]
        \begin{center}
            \begin{tabular}{|c|c|c|c|c|c|}
                \hline
                $x_1$/cm & $x_2$/cm & $x_3$/cm & $u$/cm & $v$/cm & $f$/cm \\
                \hline
                85.00 & 77.87 & 92.50 & -7.13 & 14.63 & -13.91 \\
                \hline
                85.00 & 76.60 & 99.00 & -8.40 & 22.40 & -13.44 \\
                \hline
                85.00 & 74.56 & 112.87 & -10.44 & 38.31 & -14.35 \\
                \hline
            \end{tabular}
            \caption{物像距法测量薄凹透镜焦距}
        \end{center}
    \end{table}

	\noindent \textbf{测量得到焦距平均值为：}$\bm{\bar{f} = -} \textbf{13.90\,cm}$

    \bigskip
    \subsection*{\zihao{-3} 1.4\ \ 自准直法测量薄凹透镜焦距}

    \noindent 1. 粗测凸透镜焦距$f$，取物与屏（像）间距离$A>4f$，固定屏的位置，记录屏的位置$x_1$；

    \noindent 2. 移动凸透镜，使物在屏上成小像，固定凸透镜的位置，该像将成为凹透镜的虚物，虚物的位置即为$x_1$；

    \noindent 3. 撤去屏，放置平面镜，并将凹透镜放置在凸透镜和平面镜之间，移动凹透镜和屏，使虚物在物屏上成等大倒立实像，记录此时凹透镜光心位置$x_2$与平面镜位置$x_3$；

    \noindent 3. 计算$\Delta x = -|x_2 - x_1|$，得到焦距：
    \[f = \Delta x\]

    \noindent \textbf{测量数据如下：}
    \begin{table}[ht!]
        \begin{center}
            \begin{tabular}{|c|c|c|c|c|}
                \hline
                $x_1$/cm & $x_2$/cm & $x_3$/cm & $\Delta x$/cm & $f$/cm \\
                \hline
                85.00 & 69.68 & 77.50 & -15.32 & -15.32 \\
                \hline
            \end{tabular}
            \caption{自准直法测量薄凸透镜焦距}
        \end{center}
    \end{table}

    \noindent \textbf{测量得到焦距为：}$\bm{f = -} \textbf{15.32\,cm}$

    \newpage
    \section*{\zihao{3} 2\ \ 分析与讨论}
    \bigskip

    \subsection*{\zihao{-3} 2.1\ \ 位移法和自准直法各自的优缺点}

    \noindent 1. 位移法：优点是测量较为准确，因为所测的是透镜位置的相对位移量，可以避免透镜光心与其夹具块上读数标志的位置可能不一致所导致的测量误差；缺点是操作较为繁琐。
    \medskip

    \noindent 2. 自准直法：优点是测量比较迅速，能直接测得透镜焦距的数值；缺点是误差较大，不能避免透镜光心与其夹具块上读数标志的位置可能不一致所导致的测量误差，且难以判断像的清晰度从而确定透镜的位置。

    \bigskip
    \subsection*{\zihao{-3} 2.2\ \ 实验中测量误差的来源分析}

    \noindent 1. 人眼对像清晰度的分辨能力有限，尤其是在成放大的像的情况下，光屏在一定距离范围内其上的像都足够清晰，使得对像的位置判断有一定的不确定性，像距测量误差较大。
    \medskip

    \noindent 2. 透镜光心与其夹具块上读数标志的位置可能不一致，会在自准直法等一些测量方法中导致测量误差。

    \newpage
    \section*{\zihao{3} 3\ \ 收获与感想}
    \bigskip

    \noindent 1. 位移法测量薄凸透镜焦距中，$A$越接近$4f$，大像与小像之间大小差异越小，越容易确定凸透镜的位置。若$A$过大，小像会太小以至于难以判断像的清晰度从而确定凸透镜的位置。
    \medskip

    \noindent 2. 在调节共轴时要先粗调后细调。细调共轴一次只能调节一个凸透镜或一对凸凹透镜，其过程与位移法测量薄凸透镜焦距或物像距法测量薄凹透镜焦距相同。细调薄凸透镜共轴中：成大像时移动光屏，成小像时移动凸透镜，每次移动都保证像成在光屏中心。其原理与“减半逐步逼近法”调节分光计望远镜光轴垂直于仪器转轴非常类似。
    \medskip

    \noindent 3. 物像距法测量薄凹透镜焦距中，虚物最好用凸透镜所成的小像。如果用大像，再次被凹透镜放大后所成实像的亮度过低，难以判断像的清晰度从而确定凹透镜或屏的位置。

\end{document}